\section{Limitations and Threats to Validity}
\label{sec:limitations}

\paragraph{Scope and target.}
Two assets, one window, one horizon, one convention. BTC is a single venue's USDT
price, not a global index. The label is the raw close one session ahead, so
dividends are ignored, and a one-session horizon is where persistence is
strongest --- a bias against trained models that follows from the question we
chose, and the reason H1 was pre-registered rather than discovered.

\paragraph{Statistics and evaluation integrity.}
Skill differences are within about one percentage point and intervals overlap, so
we report a tie; because errors are serially dependent, even block resampling
understates uncertainty. Every number here is a validation number, and validation
influenced model selection, so it is not a clean holdout; the sealed block exists
to give one unbiased estimate once the analysis is frozen.

\paragraph{Measurement integrity.}
Two leakage defects were found by independent audit and fixed. Both would have inflated results, and both were invisible to the authors of the
code: the audit mechanism, not the initial implementation, is what makes these
numbers trustworthy.

\paragraph{Education.}
No learner study was run, so claims about the lesson are design claims.
